\section{Lista de acciones para equipos de investigación}

\subsection{Del curso a la ingeniería: cómo empezar a construir un flujo de desarrollo de Agentes reproducible}
\begin{figure}[H]
\centering
\includegraphics[width=0.93\textwidth]{frame-08-research-questions.jpg}
\caption{Cierre de las preguntas de investigación: el triángulo de capacidad, seguridad y participación humana}
\end{figure}
\noindent\footnotesize Evidencia visual: subtítulos aproximadamente en 00:51:01--00:52:00, el docente enfatiza los temas sistémicos y las líneas de investigación posteriores.\normalsize

\begin{importantbox}{Ruta mínima recomendada para la implementación}
\begin{enumerate}
    \item Definir primero tareas verificables, y discutir después la arquitectura del Agente;
    \item Establecer primero registros y reproducción, y optimizar después la política;
    \item Hacer primero la descomposición de fallos, y buscar después el puntaje en las tablas de clasificación;
    \item Unificar primero el protocolo de evaluación, y ampliar después el trabajo paralelo del equipo.
\end{enumerate}
\end{importantbox}

\subsection{Panel recomendado de evaluación y gobernanza}
\begin{table}[H]
\centering
\begin{tabular}{p{2.8cm}p{4.1cm}p{4.6cm}}
\toprule
\textbf{Dimensión del panel} & \textbf{Métricas clave} & \textbf{Recomendación de implementación} \\
\midrule
Desempeño en la tarea & pass rate, success@k, latency & Fijar la versión de la tarea y conservar trayectorias históricas reproducibles \\
Calidad de ejecución & tasa de fallos de herramientas, tasa de reintentos, tasa de reversiones & Introducir registros de auditoría obligatorios y etiquetas de clasificación de fallos \\
Restricciones de seguridad & eventos de exceso de privilegios, alertas de operaciones sensibles & Delimitar con claridad los límites de permisos y los puntos de confirmación humana \\
Reproducibilidad científica & varianza de los resultados, consistencia del entorno, número de repeticiones del experimento & Publicar instantáneas de configuración y versiones de los scripts de evaluación \\
\bottomrule
\end{tabular}
\caption{Panel de monitoreo para que la investigación de Agentes pase de «funcionar» a «acumularse»}
\end{table}

\begin{knowledgebox}{Un diseño organizacional clave en la colaboración de equipos}
Quien dirige la investigación debe establecer la reproducibilidad como objetivo de primer orden:
los commits de código, las versiones de datos, las versiones de modelo y las versiones de evaluación deben poder formar una misma cadena de trazabilidad, para evitar «un resultado correcto pero un proceso imposible de rastrear».
\end{knowledgebox}

\begin{warningbox}{No busques la «automatización total» de manera prematura}
En tareas de alto riesgo y alta incertidumbre, una intervención humana razonable no es un retroceso, sino un mecanismo necesario para aumentar la robustez del sistema.
El verdadero objetivo es la «automatización controlada», no la «fantasía de operar sin supervisión».
\end{warningbox}

\subsection{Resumen de la sección}
En el plano de la implementación, lo más importante no es el nombre de algún marco de trabajo, sino si estas cuatro cosas pueden sostenerse a la vez: la definición de la tarea, el protocolo de evaluación, la gobernanza de registros y el control de permisos.


